% =====================================================================================================
% Section VI. Observational production functions and algorithmic progress.
% Sources: output/memos/m4_observational(.md, _review.md), m5_progress(.md, _review.md),
% m6_montecarlo(.md, _review.md); numbers verified against output/tables/*.csv (see paper/notes/writer_observational.md).
% Tables 8-9 are generated from the module CSVs (generator recorded in writer_observational.md).
% =====================================================================================================

\section{Observational Production Functions and Algorithmic Progress}\label{sec:obs}

Section~\ref{sec:wedge} took the technology as known and read labs' choices through it. This section asks the reverse
question: what do released models, used as data, reveal about the technology? Much of what is known about returns to
compute at scale, and about the pace of algorithmic progress, comes from such observational data: regressions of benchmark scores on the inputs of models released by many labs
\citep{ruan2024observational,maiapolo2024sloth,owen2024predictable,mertens2026secret}, and of perplexity on the inputs
and release dates of published models \citep{ho2024algorithmic}. These are production functions estimated from data in
which the firms chose the inputs, and \citet{whitfill2025note} and \citet{konig2026validity} warn that they may be
confounded. By Propositions~\ref{prop:transmission} and \ref{prop:dmr}, returns to compute then carry a transmission
bias whose sign depends on how labs choose compute, and time series identify the direction of biased technical change
but not its size. We find that, once the design is matched, cross-lab returns to compute agree with
experiments, whereas the split between parameters and data and the rate of algorithmic progress are weakly identified,
and part of measured progress is allocative rather than technical.

\subsection{A LaLonde-Style Benchmark}

\citet{lalonde1986evaluating} judged nonexperimental estimators by whether they reproduce an experimental benchmark. We
apply the same test to production-function estimators, taking the experiments from designed sweeps in which the
experimenter chose the inputs (Section~\ref{sec:data}). The observational sample has 128 dense base models from 38
families and 22 developers, released between March 2021 and July 2024. These models share no validation loss, so output
is the log-odds of above-chance accuracy on HellaSwag \citep{zellers2019hellaswag}, $q=\ln[\tilde a/(1-\tilde a)]$, where
$\tilde a$ is accuracy net of chance, rescaled to the unit interval. We estimate $q_i=\theta_0+\theta_N n_i+\theta_D
d_i+\omega_i+e_i$, which is Cobb--Douglas in the odds with Hicks-neutral productivity $\omega_i$, and the regression of
$q_i$ on $c_i$ alone, whose slope $\theta_C$ is what observational scaling studies report.

The comparison must respect the curvature of the technology. In the OLMo ladder and the sweep of
\citet{gadre2024language}, HellaSwag log-odds are strongly concave in $(n,d)$ ($F$-statistics of 352 and 70 for the
quadratic terms), and at $N=1$B the local data elasticity falls from 0.58 at 20 tokens per parameter to 0.09 at 500.
Released models are heavily over-trained (median $D/N$ of 213, ninetieth percentile 1,715), so they sit where an extra
token is worth little. Observational data elasticities of 0.19--0.26 thus look 27--63 percent attenuated against global
experimental values of 0.35--0.51, but much of the gap is the technology. We therefore match the design, in the spirit of
\citet{dehejia1999causal}. We fit a quadratic surface $\hat f(n,d)$ to 88 designed runs ($R^2=0.994$), keep the 57
observational models inside the convex hull of the experimental inputs, and apply each estimator twice: to the observed
outputs, and to the outputs $q^*_i=\hat f(n_i,d_i)$ that the experimental technology produces with the same inputs,
fixed effects and controls; the difference is the estimator's bias. The null is joint: a
common technology up to Hicks-neutral shifts, outputs that map one to one across evaluation harnesses, and no
transmission or selection bias.

\begin{table}[tp]
\centering
\caption{Design-Matched Comparison of Observational and Experimental Elasticities}
\label{tab:lalonde}
\footnotesize
\setlength{\tabcolsep}{2.4pt}
\begin{tabular}{@{}l ccc cc cc@{}}
\toprule
 & \multicolumn{3}{c}{Compute only, $\theta_C$} & \multicolumn{2}{c}{Parameters, $\theta_N$} & \multicolumn{2}{c}{Tokens, $\theta_D$} \\
\cmidrule(lr){2-4}\cmidrule(lr){5-6}\cmidrule(lr){7-8}
 & Obs. & Exp. & Bias & Exp. & Bias & Exp. & Bias \\
 & (1) & (2) & (3) & (4) & (5) & (6) & (7) \\
\midrule
\multicolumn{8}{@{}l}{\textit{Panel A. HellaSwag, 57 models inside the experimental support}} \\
OLS & 0.328 & 0.329 & $-$0.001 & 0.423 & 0.079 & 0.273 & $-$0.050 \\
 & (0.037) & (0.009) & (0.038) [0.99] & (0.027) & (0.091) [0.37] & (0.017) & (0.083) [0.58] \\
\addlinespace[2pt]
Year FE & 0.319 & 0.348 & $-$0.029 & 0.416 & 0.101 & 0.304 & $-$0.113 \\
 & (0.041) & (0.008) & (0.042) [0.46] & (0.027) & (0.097) [0.32] & (0.016) & (0.108) [0.30] \\
\addlinespace[2pt]
Developer FE & 0.393 & 0.370 & 0.024 & 0.448 & 0.141 & 0.293 & $-$0.090 \\
 & (0.046) & (0.009) & (0.046) [0.68] & (0.027) & (0.079) [0.12] & (0.018) & (0.166) [0.56] \\
\addlinespace[2pt]
Family FE & 0.441 & 0.409 & 0.031 & 0.445 & 0.203 & 0.348 & $-$0.262 \\
 & (0.065) & (0.012) & (0.067) [0.47] & (0.027) & (0.086) [0.05] & (0.018) & (0.162) [0.02] \\
\addlinespace[2pt]
Lab $\times$ period FE & 0.442 & 0.405 & 0.037 & 0.445 & 0.165 & 0.358 & $-$0.113 \\
 & (0.041) & (0.010) & (0.042) [0.37] & (0.027) & (0.097) [0.12] & (0.016) & (0.180) [0.44] \\
\addlinespace[2pt]
OP-style selection & 0.335 & 0.350 & $-$0.015 & 0.496 & 0.076 & 0.281 & $-$0.057 \\
 & (0.050) & (0.008) & (0.051) & (0.025) & (0.116) & (0.017) & (0.088) \\
\addlinespace[2pt]
\multicolumn{8}{@{}l}{\textit{Panel B. HellaSwag, other compute-only estimators}} \\
OLS, EIV-corrected & 0.366 & 0.329 & 0.037 &  &  &  &  \\
 & (0.042) & (0.009) & (0.043) &  &  &  &  \\
Family FE, EIV-corrected & 0.858 & 0.409 & 0.448 &  &  &  &  \\
 & (0.127) & (0.012) & (0.128) &  &  &  &  \\
Reverse regression & 0.560 & 0.352 & 0.209 &  &  &  &  \\
 & (0.126) & (0.007) & (0.126) &  &  &  &  \\
IV, FLOP/\$ ($F=0.9$) & 0.724 & 0.012 & 0.712 &  &  &  &  \\
 & (0.452) & (0.043) & (0.454) &  &  &  &  \\
\addlinespace[2pt]
\multicolumn{8}{@{}l}{\textit{Panel C. OLS, other output and surface}} \\
Composite (57) & 0.370 & 0.388 & $-$0.018 & 0.399 & 0.109 & 0.381 & $-$0.095 \\
 & (0.045) & (0.034) & (0.057) & (0.066) & (0.119) & (0.053) & (0.107) \\
\addlinespace[2pt]
Ladder-only surface (39) & 0.281 & 0.297 & $-$0.016 & 0.430 & $-$0.028 & 0.233 & $-$0.010 \\
 & (0.053) & (0.008) & (0.054) & (0.031) & (0.163) & (0.019) & (0.131) \\
\bottomrule
\end{tabular}

\begin{tablenotes}
Output is the log-odds of above-chance accuracy, $q=\ln[\tilde a/(1-\tilde a)]$ with $\tilde a=(\text{accuracy}-\text{chance})/(1-\text{chance})$. $\theta_N$ and $\theta_D$ are the coefficients on $\ln N$ and $\ln D$ in a regression of $q$ on both; $\theta_C$ is the coefficient on $\ln C=\ln 6ND$ alone. Obs.: the estimator applied to observed outputs, standard errors clustered by developer (19 clusters in Panel A). Exp.: the same estimator, with the same fixed effects and controls, applied on the same models to counterfactual outputs from the experimental technology, a quadratic surface in $(\ln N,\ln D)$ fitted to 88 designed runs (OLMo ladder with OLMo-2 7B and 13B, and Gadre et al.\ with $N\geq 0.1$B; $R^2=0.994$); standard errors from 300 bootstrap refits of the surface. Bias $=$ Obs.\ $-$ Exp.; its standard error treats the two as independent. Brackets: restricted wild-cluster bootstrap-$t$ $p$-values (Webb weights, 999 draws) for the bias, reported for the estimators and samples for which they were computed. The 57 models lie inside the convex hull of the experimental $(\ln N,\ln D)$ points (30 families, 19 developers, all with $N\le 9$B). OP-style selection: OLS with a quadratic in the estimated propensity of being an Epoch ``notable'' model. EIV: slope divided by the reliability implied by Epoch's ``Confident'' compute class. Reverse regression: inverse slope of $\ln C$ on $q$. IV: $\ln C$ instrumented by the log FLOP/s per dollar of the best data-center accelerator released at least six months before the model; $F$ is the cluster-robust first-stage statistic; in this row the Exp.\ column applies the same instrument to the experimental outputs. With developer fixed effects the same instrument has $F<0.01$ (not shown). Panel C: composite is the logit of mean chance-adjusted accuracy on HellaSwag, ARC-Challenge and Winogrande; the ladder-only surface is fitted to OLMo runs only. ARC-Challenge, Winogrande and a subsample without imputed-data, distilled, synthetic-data and code models are in Online Appendix Table~\ref{tab:app-lalonde-robust}.
\end{tablenotes}
\begin{tablenotes}[Source]
Authors' calculations. Observational outputs: Open LLM Leaderboard v1 via \citet{ruan2024observational} and \citet{maiapolo2024sloth}; parameter counts from Hugging Face; release dates and compute from \citet{epochai2026data}. Experimental runs: \citet{bhagia2024establishing}, \citet{teamolmo2024olmo} and \citet{gadre2024language}.
\end{tablenotes}
\end{table}


\begin{figure}[tp]
\centering
\includegraphics[width=\textwidth]{m4_observational_fig5_lalonde}
\caption{Observational Estimates and the Experimental Technology on the Same Design}\label{fig:lalonde}
\begin{figurenotes}
Circles: estimates from the observed HellaSwag log-odds of the 57 released models inside the experimental support,
with 95 percent confidence intervals clustered by developer. Diamonds: the same estimator applied to the outputs that the
experimental technology produces for the same models, with 95 percent intervals from 300 bootstrap refits of the
surface. Shaded bands: range of global log-linear fits in the designed sweeps (not design-matched). Instrumental-variable
rows are omitted (first-stage $F<1$). Estimates are in Table~\ref{tab:lalonde}.
\end{figurenotes}
\end{figure}

Returns to compute show no detectable bias (Table~\ref{tab:lalonde}, Figure~\ref{fig:lalonde}). Pooled OLS gives
$\theta_C=0.328$ (0.037) against 0.329 (0.009) from the experimental technology on the same models, a bias of
$-0.001$ (0.038) with a wild-cluster bootstrap $p$-value of 0.99 \citep{cameron2008bootstrap,webb2023reworking}. Across
six estimators, from pooled OLS to lab-by-period fixed effects, the biases lie between $-0.03$ and $+0.04$, and none is
significant; a three-benchmark composite and a surface fitted to the OLMo ladder alone agree (panel C). This is a
failure to
reject, not proof of unbiasedness. The 95 percent interval for the OLS bias, $[-0.075,0.074]$, spans about $\pm23$
percent of the benchmark; point estimates for ARC-Challenge and Winogrande, although insignificant, would imply 18--31
percent attenuation; and a subsample without code, distilled, synthetic-data and imputed-data models has a bias of 0.045
(0.023), with wild $p=0.10$ (Online Appendix Table~\ref{tab:app-lalonde-robust}). The test covers only models with at most 9B parameters, and its benchmark is the technology
of the AI2 and OpenLM recipes.

The split between parameters and data is less reliable. With family fixed effects, the parameter elasticity is
overstated by 0.20 (0.09; wild $p=0.05$) and the data elasticity understated by 0.26 (0.16; wild $p=0.02$), but the
pattern is fragile: pooled OLS biases are insignificant, ARC-Challenge and Winogrande give negative biases for
$\theta_N$, and the ladder-only surface reverses the signs under developer fixed effects. Within families the standard
deviation of $\ln D$ is 0.42, against 1.29 for $\ln N$, and only 18 of 38 families vary $D$ at all. Moreover,
transmission of Hicks-neutral productivity loads on $c=\ln6+n+d$, so it would bias both elasticities in the same
direction (Proposition~\ref{prop:transmission}). Biases of opposite sign point instead to something correlated with the
input mix, such as measurement error in $D$ or data-quality changes within a family, which this benchmark cannot
distinguish from bias. Both worlds agree that $\theta_N>\theta_D$ on this design (0.42 against 0.27 under the
experimental technology): over-trained models have wedges above one (Section~\ref{sec:wedge}), so regressing on
compute alone imposes a restriction, $\theta_N=\theta_D$, that the data reject ($\theta_N-\theta_D=0.21$, standard error
0.08, on all 128 models).

How large would transmission bias be under different compute rules? Treating the 25 pretraining recipes of
\citet{magnusson2025datadecide} as labs with real outcomes and simulating how they choose model size, pooled OLS is
unbiased under exogenous budgets, overstates $\theta_C$ by 0.044 when funding rises with productivity rank, and
understates it by 0.18 (40 percent) when each lab trains to a common capability target. Cross-family productivity
dispersion in the observational panel is about five times that of the recipes, which scales the funding bias to about
$+0.2$ and the target bias beyond $-0.18$. The observed intervals exclude both. The data are consistent with budgets
roughly orthogonal to productivity, or with offsetting funding and release selection; the simulated rules are stylized.

The standard IO remedies cannot be used on these data, for reasons that concern the data rather than the estimators.
The price-performance of the best available accelerator is a step function of calendar time, with first-stage $F$ at
most 3.8 in every specification; that of each model's own hardware ($F=10.7$) takes three values identified by six
models, effectively a vintage dummy. An export-control instrument (Chinese developers after October 2023) is strong among models below 14B
parameters ($F=10.6$) but gives $\theta_C=0.90$ (0.25), 1.7--2.3 times the OLS and fixed-effects estimates on that sample;
its exclusion restriction fails because recipes and data changed at the same time. Dynamic-panel estimators \citep{arellano1991some,blundell1998initial} need three consecutive
generations, which only six product-line cells have. The proxy inversion of \citet{olley1996dynamics} and
\citet{levinsohn2003estimating} is unavailable by construction, because Hicks-neutral productivity does not move the
input mix at given compute (Online Appendix A, Proposition~\ref{prop:proxy}). Measurement error in compute is not the
problem: where Epoch reports independent compute estimates, the implied reliability of $\ln C$ is at least 0.988.

The panel also measures productivity dispersion (Online Appendix Table~\ref{tab:app-tfp}). In compute-equivalent units, the 90--10 ratio of family effects is
23.7 (90 percent interval $[9.1,197]$), and 7--9 without code, distilled and synthetic-data families; within developers,
as in \citet{mertens2026secret}, it is 11.1 $[4.4,18.8]$, against their 41 on another benchmark. Since returns to scale
are near one, the manufacturing ratio of 1.92 \citep{syverson2011determines} is comparable, and LLM dispersion is three
to five times larger in logs. The synthetic-data families Phi and SmolLM rank first and third of 38: training compute
omits the teacher's compute, so measured productivity credits the student with an intermediate input, the gross-output
problem of \citet{gandhi2020identification}.

What data would make the remedies work? Online Appendix C (Table~\ref{tab:app-mc-industry}) simulates an industry of labs that release one to four models
per generation, with the \citet{besiroglu2024chinchilla} technology, persistent and growing Hicks-neutral productivity,
heterogeneous inference demand and four compute rules. Pooled estimation overstates the frontier elasticity $\gamma$ by
13 percent when funding responds to productivity, understates it by 29 percent under a capability target, and
overstates it by 66 percent with one flagship per lab and generation, as Proposition~\ref{prop:transmission} predicts.
Estimation within lab and generation is the observational analogue of an IsoFLOP experiment: siblings share productivity
and differ in size for reasons unrelated to it, so the estimator is unbiased under every rule, but it needs at least
two released siblings, and ACF-type GMM succeeds only with the same within-family variation. A first-order-condition system
that imposes training-only optimality ($w=1$) overstates $\ln M^*$ by 1.25--1.28, exactly $2E[\ln w]/(\alpha+\beta)$.
Pooled fits with $E$ free and no time effects, as in ML practice, attribute 54--60 percent of productivity growth to the
scaling law. The calibration is assumed, so only signs and rankings carry over.

\subsection{Algorithmic Progress}

\citet{ho2024algorithmic} estimate that the compute needed to reach a given perplexity halves every 8.4 months, from 231
model--benchmark observations in 144 papers. Their preferred model is the technology of Section~\ref{sec:framework} with
$E=0$ and factor-augmenting productivity growing linearly in time, $\psi_N=g_Nt$ and $\psi_D=g_Dt$. Effective compute
grows at $g_C=g_N+g_D$ and doubles every $T_C=12\ln2/g_C$ months; Hicks neutrality is $\alpha g_N=\beta g_D$. Because the
sample is not on the expansion path, a problem of the kind studied by \citet{diamond1978measurement} arises. To first
order, the data pin down how fast loss falls at fixed inputs, an average of $\alpha g_N$ and $\beta g_D$ weighted by the
mean shares of the two terms in loss (0.55 and 0.45), whereas the doubling time weights the two rates by $1/\alpha$ and
$1/\beta$. The two combinations coincide only on the expansion path, where the parameter share would be $a=0.37$ under
their exponents; off it, rate pairs that fit equally well imply different doubling times.

\begin{table}[tp]
\centering
\caption{Algorithmic Progress under Alternative Estimators and Identifying Assumptions}
\label{tab:progress}
\footnotesize
\setlength{\tabcolsep}{3pt}
\begin{tabular}{@{}l ccccccc@{}}
\toprule
 & $\hat\alpha$ & $\hat\beta$ & \multicolumn{1}{c}{$g_N$} & \multicolumn{1}{c}{$g_D$} & $T_C$ & 95\% CI & MSE \\
\midrule
\multicolumn{8}{@{}l}{\textit{Panel A. Doubling time of effective compute, data of Ho et al.\ (2024)}} \\
Published (bootstrap median) & 0.068 & 0.040 &  &  & 8.4 & $[4.5, 14.3]$ & -- \\
(1) Replication, default stop & 0.068 & 0.040 & 0.06 & 0.90 & 8.7 & $[4.1, 16.9]$ & 0.0464 \\
(2) Run to convergence & 0.138 & 0.025 & $-$0.43 & 1.80 & 6.1 & $[3.0, 22.7]$ & 0.0463 \\
(3) Unpenalized NLS & 0.078 & 0.031 & 0.35 & 0.46 & 10.2 & $[4.1, 27.6]$ & 0.0453 \\
\quad profile-likelihood CI & & & & & & $[4.1, 40.5]$ & \\
(4) $E_b$ estimated & 0.072 & 0.042 & 0.34 & 0.45 & 10.5 & $[2.9, 26.3]$ & 0.0451 \\
(5) Hicks-neutral, as (2) & 0.059 & 0.050 & 0.33 & 0.39 & 11.6 & $[6.6, 17.6]$ & 0.0465 \\
(6) Hicks-neutral, as (3) & 0.083 & 0.030 & 0.24 & 0.66 & 9.2 & $[6.3, 14.3]$ & 0.0454 \\
(7) Effective data & 0.057 & 0.067 & 0.58 & $-$0.25 & 25.7 & $[3.9, \infty)$ & 0.0484 \\
(8) Impose $\gamma=0.0525$ & 0.058 & 0.574 & 0.22 & 0.17 & 21.5 & $[9.9, 48.2]$ & 0.0509 \\
(9) Year terms $\div$ exponents & 0.348 & 0.366 & 0.01 & 0.10 & 76.5 & $[46.2, \infty)$ & -- \\
(10) Besiroglu $(\alpha,\beta)$ & 0.348 & 0.366 & $-$0.39 & 0.28 & $g_C<0$ & $[72.5, \infty)$ & 0.1126 \\
(11) Hoffmann $(\alpha,\beta)$ & 0.339 & 0.285 & $-$0.44 & 0.35 & $g_C<0$ & $[54.3, \infty)$ & 0.0994 \\
(12) 12 sweep technologies &  &  &  &  & \multicolumn{2}{c}{$g_C<0$ or 42--122} & 0.060--0.219 \\
\addlinespace[3pt]
\multicolumn{8}{@{}l}{\textit{Panel B. Allocative gain, models with $C\geq10^{23}$ FLOP, 2020--21 to 2022--24}} \\
Technology & $n_1,n_2$ & $\overline{CE}_1$ & \multicolumn{1}{c}{$\overline{CE}_2$} & \multicolumn{1}{c}{Gain} & 95\% CI & & Share \\
\cmidrule(l){1-8}
Besiroglu et al. & 7, 92 & 0.57 & 0.60 & 1.05 & $[0.57, 1.72]$ &  & 0.02 \\
Hoffmann et al., TeX & 7, 92 & 0.38 & 0.73 & 1.91 & $[1.49, 2.43]$ &  & 0.30 \\
Hoffmann et al., rounded & 7, 92 & 0.31 & 0.74 & 2.41 & $[1.92, 3.14]$ &  & 0.40 \\
Farseer grid & 7, 92 & 0.44 & 0.70 & 1.58 & $[1.27, 2.04]$ &  & 0.21 \\
Nine sweep technologies & 7, 92 &  &  & 0.55--4.31 &  &  & $-$0.27 to 0.67 \\
Besiroglu, cleaned & 7, 79 & 0.57 & 0.62 & 1.09 & $[0.58, 1.74]$ &  & 0.04 \\
Besiroglu, top five & 10, 15 & 0.67 & 0.85 & 1.27 & $[0.96, 1.77]$ &  & 0.10 \\
Ho et al.\ (2024) & 7, 92 & 0.10 & 0.01 & 0.09 & $[0.07, 0.13]$ &  & $-$1.10 \\
\bottomrule
\end{tabular}

\begin{tablenotes}
Panel A re-estimates the preferred model of \citet{ho2024algorithmic}, $\ln\text{ppl}=\exp(\alpha_{0b}-\alpha g_N t-\alpha\ln N)+\exp(\beta_{0b}-\beta g_D t-\beta\ln D)$ with irreducible loss set to zero, under the stated estimator or restriction. $g_N,g_D$: parameter- and data-augmenting growth rates (log points per year); $T_C=12\ln2/(g_N+g_D)$: effective-compute doubling time in months. Rows (1)--(2) minimize their objective, mean squared error plus $0.0025\sum|\theta|$, by SLSQP from zero; row (1) uses the scipy default tolerance of the authors' code and row (2) runs to convergence. The published 8.4 months is a bootstrap median; the point estimate of row (1) is 8.7. Intervals: 95\% paper-cluster bootstrap (400 draws), percentiles of $g_N+g_D$ mapped to $T_C$ ($\infty$ when non-positive growth lies inside the interval); the published interval is the authors' iid bootstrap; the profile-likelihood interval inverts an iid Gaussian likelihood-ratio test. Hicks neutrality: $\alpha g_N=\beta g_D$. (7): data measured as effective data under repetition \citep{muennighoff2023scaling}. (8): the frontier elasticity obtained by fitting the $E=0$ form to the Chinchilla sweep. (9): Ho et al.'s year coefficients divided by the exponents of \citet{besiroglu2024chinchilla}, without re-estimation \citep{whitfill2025note}. (10)--(12): exponents fixed at experimental values, benchmark-specific $E_b$ estimated; row (12) reports the range over the 12 technologies estimated in Section~\ref{sec:tech} (two on the Chinchilla sweep): $g_C<0$ for 9 of them; for the other 3, $T_C$ lies between 42 and 122 months, and every interval includes non-positive growth. Panel B: $\overline{CE}$ is the geometric mean of Farrell cost efficiency $C_{\min}(L(N,D))/6ND$ under the stated technology in each period ($n_1$, $n_2$ models); gain $=\overline{CE}_2/\overline{CE}_1$, with a bootstrap over models (and over parameter draws for Besiroglu et al.). Share $=\ln(\text{gain})/\ln(8.9)$, where 8.9 is the effective-compute gain implied by row (1) over the difference in mean release dates (9.9 for the top-five sample, 8.2 for the cleaned sample); a negative share means that cost efficiency fell. Top five: the five models with the most compute in each year. Cleaned: drops derivatives, systems-paper runs, non-transformers, exact re-releases, models with $D/N$ exactly 20 and Epoch ``speculative'' entries.
\end{tablenotes}
\begin{tablenotes}[Source]
Authors' calculations from the data and code of \citet{ho2024algorithmic} (Panel A) and the Epoch AI models database \citep{epochai2026data} (Panel B).
\end{tablenotes}
\end{table}


We reproduce Ho et al.'s estimates exactly (Table~\ref{tab:progress}; Online Appendix
Table~\ref{tab:app-ho-replication}): the same objective value, all parameters within
$4\times10^{-7}$, and a bootstrap median of 8.44 months under their protocol. The published 8.4 months is that median;
the point estimate is 8.7. It is not the minimizer of the authors' objective. SciPy's default stopping rule halts the
optimizer after 18 iterations, part of the way along the ridge. Run to convergence, the objective falls from 0.0518 to
0.0507, all 300 random starts end below the published value, and $T_C=6.1$ months with a paper-cluster interval of
$[3.0,22.7]$, twice as wide as published. We do not claim that 6.1 months is right: mean squared error changes by only
0.27 percent, the gap is almost entirely an L1 penalty on constants that depend on an arbitrary normalization, and
unpenalized least squares gives 10.2 months $[4.1,27.6]$. Estimators of the same model thus give 6.1--10.2 months, and
the data cannot arbitrate between them. This is the failure that \citet{besiroglu2024chinchilla} found in the third approach of
\citet{hoffmann2022training}, and it echoes convergence problems in demand estimation
\citep{knittel2014estimation,dube2012improving}.

\begin{figure}[tp]
\centering
\includegraphics[width=\textwidth]{m5_progress_dmr_ridge}
\caption{The Ridge in Estimates of Algorithmic Progress}\label{fig:ridge}
\begin{figurenotes}
Data of \citet{ho2024algorithmic}, 231 model--benchmark observations; $\alpha_{\text{year}}\equiv\alpha g_N$ and
$\beta_{\text{year}}\equiv\beta g_D$. Panel (a): 95 percent joint profile-likelihood region, bootstrap draws from the
authors' code, the optimizer's iterates from its default stopping point to the converged minimum of the same objective,
the unpenalized least-squares minimum, and the Hicks-neutral line. Panel (b): effective-compute doubling time $T_C$
along the floor of the profile inside the 95 percent region. Panel (c): profile likelihood-ratio statistic over the
parameter-augmenting share $\phi=g_N/g_C$, with the implied $T_C$ in months; the dashed line is the 95 percent cutoff.
Likelihood ratios assume independent Gaussian errors.
\end{figurenotes}
\end{figure}

Figure~\ref{fig:ridge} shows the ridge. The bootstrap correlation between $\alpha g_N$ and $\beta g_D$ is $-0.73$ to
$-0.82$. The profile-likelihood interval for $T_C$ is $[4.1,40.5]$ months (a paper-cluster bootstrap gives
$[4.1,27.6]$), and the parameter-augmenting share $\phi=g_N/g_C$ can lie anywhere in $[-0.39,3.37]$, over which $T_C$
ranges from 5 to 38 months. Hicks neutrality is not rejected (bootstrap $p=0.29$ and 0.94 for the converged and unpenalized estimators), and
imposing it gives 9--12 months with the tightest intervals among specifications the data do not reject. As
\citet{diamond1978measurement} predict, precision about the rate is bought with an assumption about the bias, and even
that assumption is fragile: if the true technology has $E>0$, neutrality imposed within the $E=0$ form is misspecified,
and in simulations it overstates progress.

Ho et al.'s data exponent is 0.04, against about 0.37 in experiments, which \citet{whitfill2025note} reads as progress
overstated about ninefold. Imposing experimental exponents is rejected (mean squared error rises by a factor of 1.3--4.8
across 14 technologies) and implies negative or implausibly slow progress. Two specification choices explain much of the
exponent gap while moving $T_C$ far less (Online Appendix Table~\ref{tab:app-attenuation}). First, fitting the $E=0$ form to the Chinchilla sweep itself lowers the
frontier elasticity from 0.178 to 0.0525, close to the sweep's mean elasticity of total rather than reducible loss
(0.054), so the like-for-like gap is 2--2.5-fold, not 7--9-fold. Dropping $E$ rescales exponents and year coefficients
by similar factors, so rates barely move: imposing a common $E$ from 0 to 1.5 nats per word doubles $\hat\gamma$ but
moves $T_C$ only from 10.2 to 12.4 months. (Ho et al.'s data themselves prefer $E=0$ within this form, so the argument
rests on the sweep and on simulations.) Second, their $D$ is dataset size, while small models were trained for a median
of 62--140 epochs; measuring effective data under repetition \citep{muennighoff2023scaling} raises $T_C$ by a factor of
1.5--4, with an unbounded upper confidence limit. Finally, frontier regressions that omit time effects overstate the compute
elasticity by a factor of 1.15 $[1.10,1.21]$ on all models and 1.44 $[1.24,1.78]$ on record-setting models, the Sahal bias (Online Appendix A,
Corollary~\ref{cor:sahal}, and Table~\ref{tab:app-ceg-sahal}; \citealp{nagy2013statistical}).

\subsection{Allocative versus Technical Change}

Measured progress, the fall in the compute needed to reach a given loss, mixes shifts of the frontier with movements
toward it. The Farrell decomposition (Online Appendix A, Proposition~\ref{prop:farrell}) splits a model's log excess
compute into an allocative term that depends only on its wedge (Section~\ref{sec:framework}) and a technical term. The
move from the allocations of \citet{kaplan2020scaling} to those of \citet{hoffmann2022training} changed the mix, not the
technology, so its gain is allocative; under the \citet{besiroglu2024chinchilla} parameters, Gopher's allocation used
about twice the compute its loss required. \citet{gundlach2025origin} count such rebalancing as algorithmic progress.

Panel B of Table~\ref{tab:progress} (and Online Appendix Figure~\ref{fig:app-allocative}) measures the realized gain for
models in the Epoch database with at least $10^{23}$ FLOP: 7 released in 2020--21 and 92 in 2022--24. Under the \citet{besiroglu2024chinchilla} technology, geometric-mean
cost efficiency rose from 0.57 to 0.60, a gain of 1.05 $[0.57,1.72]$; the gain is 1.9--2.4 under the two versions of
the \citet{hoffmann2022training} parameters and 1.6 under the Farseer grid. The median ratio of tokens to parameters rose
from 1.7 to 92: the correction of under-training was offset by deliberate over-training, which cost efficiency measured
against the training-only frontier counts as waste (Section~\ref{sec:wedge}). Against the 8.9-fold gain implied by Ho et
al.'s point estimate over the same interval, the allocative share is 2--40 percent (and the denominator ranges from 6.4 to 22.5
across the estimators of Panel A). The answer depends on the technology: across nine technologies estimated in
Section~\ref{sec:tech} the gain ranges from 0.55 to 4.3, and Meta's Llama 3 IsoFLOP runs give 0.74, so not even the sign
is robust. The defensible statement is that the realized allocative gain was at most comparable to, and under the
Chinchilla-family and Farseer technologies a small fraction of, the gain implied by Ho et al.'s rate; with seven
early models, the comparison is imprecise.

The 10-fold gain that \citet{gundlach2025origin} attribute to rebalancing is a counterfactual. Continuing Kaplan's rule, anchored at
GPT-3, would cost a factor of 1.6 at GPT-3's compute, 5.2 at $3.8\times10^{25}$ FLOP and 11.7 at $5\times10^{26}$ FLOP
under the \citet{besiroglu2024chinchilla} technology (2.4, 11.5 and 32 under Hoffmann et al.'s). The 10-fold figure is
thus the cost that Kaplan's rule would impose at 2025 frontier compute, not a gain realized between eras. Under Ho et
al.'s own exponents, by contrast, cost efficiency falls tenfold between eras: an attenuated data exponent makes tokens
nearly worthless, so a regression with those exponents cannot credit rebalancing to inputs and loads it on the time trend.

Observational data thus deliver returns to compute, once the design is matched, and the dispersion of productivity across labs. They
do not deliver the split between inputs, the bias of technical change, or its rate without a restriction on the bias.
Proposition~\ref{prop:dmr}
names the missing data: IsoFLOP sweeps at several vintages on a common output, whose drift in the compute-optimal $M^*$
signs the bias and whose curvature fixes its size. \textcolor{red}{[TBD-m9: one sentence on what our two-corpus
experiment (Section~\ref{sec:tech}) implies for the sign of the bias of a data-quality improvement, e.g., the shift in
compute-optimal $M^*$ at fixed compute.]}
